\documentclass[12pt,letterpaper]{article}
\usepackage{graphicx,textcomp}
\usepackage{natbib}
\usepackage{setspace}
\usepackage{fullpage}
\usepackage{color}
\usepackage[reqno]{amsmath}
\usepackage{amsthm}
\usepackage{fancyvrb}
\usepackage{amssymb,enumerate}
\usepackage[all]{xy}
\usepackage{endnotes}
\usepackage{lscape}
\newtheorem{com}{Comment}
\usepackage{float}
\usepackage{hyperref}
\newtheorem{lem} {Lemma}
\newtheorem{prop}{Proposition}
\newtheorem{thm}{Theorem}
\newtheorem{defn}{Definition}
\newtheorem{cor}{Corollary}
\newtheorem{obs}{Observation}
\usepackage[compact]{titlesec}
\usepackage{dcolumn}
\usepackage{tikz}
\usetikzlibrary{arrows}
\usepackage{multirow}
\usepackage{xcolor}
\newcolumntype{.}{D{.}{.}{-1}}
\newcolumntype{d}[1]{D{.}{.}{#1}}
\definecolor{light-gray}{gray}{0.65}
\usepackage{url}
\usepackage{listings}
\usepackage{color}

\definecolor{codegreen}{rgb}{0,0.6,0}
\definecolor{codegray}{rgb}{0.5,0.5,0.5}
\definecolor{codepurple}{rgb}{0.58,0,0.82}
\definecolor{backcolour}{rgb}{0.95,0.95,0.92}

\lstdefinestyle{mystyle}{
	backgroundcolor=\color{backcolour},   
	commentstyle=\color{codegreen},
	keywordstyle=\color{magenta},
	numberstyle=\tiny\color{codegray},
	stringstyle=\color{codepurple},
	basicstyle=\footnotesize,
	breakatwhitespace=false,         
	breaklines=true,                 
	captionpos=b,                    
	keepspaces=true,                 
	numbers=left,                    
	numbersep=5pt,                  
	showspaces=false,                
	showstringspaces=false,
	showtabs=false,                  
	tabsize=2
}
\lstset{style=mystyle}
\newcommand{\Sref}[1]{Section~\ref{#1}}
\newtheorem{hyp}{Hypothesis}


\title{Problem Set 4}
\date{Due: November 26, 2026}
\author{Applied Stats/Quant Methods 1}


\begin{document}
	\maketitle
	\section*{Instructions}
	\begin{itemize}
		\item Please show your work! You may lose points by simply writing in the answer. If the problem requires you to execute commands in \texttt{R}, please include the code you used to get your answers. Please also include the \texttt{.R} file that contains your code. If you are not sure if work needs to be shown for a particular problem, please ask.
		\item Your homework should be submitted electronically on GitHub.
		\item This problem set is due before 23:59 on Thursday November 26, 2026. No late assignments will be accepted.
	\end{itemize}



	\vspace{.5cm}
\section*{Question 1: Economics}
\vspace{.25cm}
\noindent 	
In this question, use the \texttt{prestige} dataset in the \texttt{car} library. First, run the following commands:

\begin{verbatim}
install.packages(car)
library(car)
data(Prestige)
help(Prestige)
\end{verbatim} 


\noindent We would like to study whether individuals with higher levels of income have more prestigious jobs. Moreover, we would like to study whether professionals have more prestigious jobs than blue and white collar workers.

\newpage
\begin{enumerate}
	
	\item [(a)]
	Create a new variable \texttt{professional} by recoding the variable \texttt{type} so that professionals are coded as $1$, and blue and white collar workers are coded as $0$ (Hint: \texttt{ifelse}).
	
	\vspace{.5cm}
	\lstinputlisting[language=R, firstline=42, lastline=42]{PS04_answer_ED.R}
	
	\item [(b)]
	Run a linear model with \texttt{prestige} as an outcome and \texttt{income}, \texttt{professional}, and the interaction of the two as predictors (Note: this is a continuous $\times$ dummy interaction.)
	
	\vspace{.5cm}
	\lstinputlisting[language=R, firstline=46, lastline=47]{PS04_answer_ED.R}
	% Table created by stargazer v.5.2.3 by Marek Hlavac, Social Policy Institute. E-mail: marek.hlavac at gmail.com
	% Date and time: Thu, Nov 26, 2026 - 12:14:08
	\begin{table}[!htbp] \centering 
		\caption{} 
		\label{} 
		\begin{tabular}{@{\extracolsep{5pt}}lc} 
			\\[-1.8ex]\hline 
			\hline \\[-1.8ex] 
			& \multicolumn{1}{c}{\textit{Dependent variable:}} \\ 
			\cline{2-2} 
			\\[-1.8ex] & prestige \\ 
			\hline \\[-1.8ex] 
			income & 0.003$^{***}$ \\ 
			& (0.0005) \\ 
			& \\ 
			professional & 37.781$^{***}$ \\ 
			& (4.248) \\ 
			& \\ 
			professional & $-$0.002$^{***}$ \\ 
			& (0.001) \\ 
			& \\ 
			Constant & 21.142$^{***}$ \\ 
			& (2.804) \\ 
			& \\ 
			\hline \\[-1.8ex] 
			Observations & 98 \\ 
			R$^{2}$ & 0.787 \\ 
			Adjusted R$^{2}$ & 0.780 \\ 
			Residual Std. Error & 8.012 (df = 94) \\ 
			F Statistic & 115.878$^{***}$ (df = 3; 94) \\ 
			\hline 
			\hline \\[-1.8ex] 
			\textit{Note:}  & \multicolumn{1}{r}{$^{*}$p$<$0.1; $^{**}$p$<$0.05; $^{***}$p$<$0.01} \\ 
		\end{tabular} 
	\end{table} 
	
	\item [(c)]
	Write the prediction equation based on the result.
	
	First I have noted the formula for the prediction equation, and, below it, the prediction equation with the values that we obtained from the regression model added:
	
	$\makebox[\linewidth]{\^{y} = $\beta_0$ + $\beta_1$ * income + $\beta_2$ * professional + $\beta_3$(income * professional)}$
	
	$\makebox[\linewidth]{\^{y}  = 21.142 + 0.0031* income + 37.7812*professional - 0.0023*income*professional}$
	
	\vspace{.5cm}
	More precisely, we can now formulate the prediction equation for professionals and for non-professionals, by inputting 1 for professionals, and 0 for non-professionals. The equations are as follows:
	
	$\makebox[\linewidth]{\^{y} = 21.142 + 0.0031* income + 37.7812*1 - 0.0023*income*1}$
	
	$\makebox[\linewidth]{\^{y}  = 21.142 + 0.0031* income}$
	\vspace{.5cm}
	\item [(d)]
	Interpret the coefficient for \texttt{income}.
	
	\vspace{.5cm}	
	For non-professionals, a one unit increase in income is associated with an increase of 0.0031 in prestige of their job, on average.
	\vspace{.5cm}
	\item [(e)]
	Interpret the coefficient for \texttt{professional}.
	\vspace{.5cm}
	At an income of zero, professionals have, on average, a higher prestige value of 37.78 for their job than non-professionals, i.e. blue and white collar workers. The interaction effect of professional status on income decreases the prestige by -0.0023.
	
	\vspace{.5cm}
	\item [(f)]
	What is the effect of a \$1,000 increase in income on prestige score for professional occupations? In other words, we are interested in the marginal effect of income when the variable \texttt{professional} takes the value of $1$. Calculate the change in $\hat{y}$ associated with a \$1,000 increase in income based on your answer for (c).
	
	\vspace{.5cm}
	
	The effect can be obtained by inputting an income value into the prediction equation for professionals.
	
\begin{verbatim}
	> # for 0$
	> 21.142 + 0.0031 * 0 + 37.7812 * 1 - 0.0023*0*1
	[1] 58.9232
	> #result = 58.9232
	> # for 1,000$
	> 21.142 + 0.0031* 1000 + 37.7812*1 - 0.0023*1000*1
	[1] 59.7232
	> # Marginal change
	> 59.7232 - 58.9232
	[1] 0.8
	> # Check - marginal change of X1 on Y: ß1 + ß3 * X2
	> 0.0031 + (-0.0023) * 1
	[1] 8e-04
\end{verbatim}
	The marginal change of income on prestige is 0.0008 at an income of $1$. For an income of $1,000$, the change in prestige is 0.8. 
	\vspace{.5cm}
	\item [(g)]
	What is the effect of changing one's occupations from non-professional to professional when her income is \$6,000? We are interested in the marginal effect of professional jobs when the variable \texttt{income} takes the value of $6,000$. Calculate the change in $\hat{y}$ based on your answer for (c).
	\vspace{.5cm}
	The calculation can be done with the prediction equation, inserting the income value of $6,000$ for both non-professionals and for professionals. This results in the following calculation and prestige values:
	
	\begin{verbatim}
		> # Non-professionals:
		> 21.142 + 0.0031* 6000 + 37.7812*0 - 0.00232*6000*0
		[1] 39.742
		> # result = 39.742
		> # Professionals:
		> 21.142 + 0.0031* 6000 + 37.7812*1 - 0.00232*6000*1
		[1] 63.6032
	\end{verbatim}
	Subtracting the results from one another, we obtain the change in prestige between a non-professional and a professional with the same income of $6,000$:
	
\begin{verbatim}
	> 63.6032 - 39.742
	[1] 23.8612
\end{verbatim}

A change from non-professional to professional status at an income of $6,000$ is associated with an increase in prestige of 23.8612.
	
\end{enumerate}

\newpage

\section*{Question 2: Political Science}
\vspace{.25cm}
\noindent 	Researchers are interested in learning the effect of all of those yard signs on voting preferences.\footnote{Donald P. Green, Jonathan	S. Krasno, Alexander Coppock, Benjamin D. Farrer,	Brandon Lenoir, Joshua N. Zingher. 2016. ``The effects of lawn signs on vote outcomes: Results from four randomized field experiments.'' Electoral Studies 41: 143-150. } Working with a campaign in Fairfax County, Virginia, 131 precincts were randomly divided into a treatment and control group. In 30 precincts, signs were posted around the precinct that read, ``For Sale: Terry McAuliffe. Don't Sellout Virgina on November 5.'' \\

Below is the result of a regression with two variables and a constant.  The dependent variable is the proportion of the vote that went to McAuliff's opponent Ken Cuccinelli. The first variable indicates whether a precinct was randomly assigned to have the sign against McAuliffe posted. The second variable indicates
a precinct that was adjacent to a precinct in the treatment group (since people in those precincts might be exposed to the signs).  \\

\vspace{.5cm}
\begin{table}[!htbp]
	\centering 
	\textbf{Impact of lawn signs on vote share}\\
	\begin{tabular}{@{\extracolsep{5pt}}lccc} 
		\\[-1.8ex] 
		\hline \\[-1.8ex]
		Precinct assigned lawn signs  (n=30)  & 0.042\\
		& (0.016) \\
		Precinct adjacent to lawn signs (n=76) & 0.042 \\
		&  (0.013) \\
		Constant  & 0.302\\
		& (0.011)
		\\
		\hline \\
	\end{tabular}\\
	\footnotesize{\textit{Notes:} $R^2$=0.094, N=131}
\end{table}

\vspace{.5cm}
\begin{enumerate}
	\item [(a)] Use the results from a linear regression to determine whether having these yard signs in a precinct affects vote share (e.g., conduct a hypothesis test with $\alpha = .05$).
	
	We perform a t-test with the coefficient for precincts with yard signs where our null hypothesis is H0: $\beta_1$ = 0, and our alternative hypothesis is H1: $\beta_1$ $\neq$ 0.
	
	First, we calculate the t-statistic and the degrees of freedom.
	
\begin{verbatim}
	> # t = beta1_hat / standard error
	> tstat1 <- 0.042 / 0.016
	> # t = beta1_hat / standard error
	> tstat1 <- 0.042 / 0.016
	> tstat1
	[1] 2.625
	> # degrees of freedom = n - k - 1
	> df <- 131 - 2 - 1
	> # degrees of freedom = n - k - 1
	> df <- 131 - 2 - 1
	> df
	[1] 128
\end{verbatim}

	
	With the t-statistic and the degrees of freedom we can calculate the p-value.
	
	\lstinputlisting[language=R, firstline=126, lastline=126]{PS04_answer_ED.R}
	
	\begin{verbatim}
		> 2*pt(tstat1, df, lower.tail=FALSE)
		[1] 0.00972002
	\end{verbatim}
	
	Based on the p-value of 0.0097, we can reject the null hypothesis that the presence of yard signs in a precinct does not affect vote share at a significant level of 0.05.
	
	\vspace{.5cm}
	\item [(b)]  Use the results to determine whether being
	next to precincts with these yard signs affects vote
	share (e.g., conduct a hypothesis test with $\alpha = .05$).
	
	A hypothesis test can be conducted in the same fashion as above, using the coefficient for the precincts next to those with yard signs this time.
	
	H0: $\beta_2$ = 0; H1: $\beta_2$ $\neq$ 0
	
\begin{verbatim}
	> tstat2 <- 0.042 / 0.013
	> tstat2
	[1] 3.230769
\end{verbatim}

The degrees of freedom stay the same, and we calculate the p-value again.

\lstinputlisting[language=R, firstline=142, lastline=142]{PS04_answer_ED.R}

\begin{verbatim}
	> 2*pt(tstat2, df, lower.tail=FALSE)
	[1] 0.00156946
\end{verbatim}

Based on  the p-value of 0.0015, we can reject the null hypothesis that being next to precincts with yard signs has no effect on vote share at a significance level of 0.05.

	\vspace{7cm}
	\item [(c)] Interpret the coefficient for the constant term substantively.
	\vspace{.5cm}
	
	In precincts without proximity to lawn signs, the vote share of Ken Cuccinelli is 30.2\% on average.
	\vspace{.5cm}
	\item [(d)] Evaluate the model fit for this regression.  What does this	tell us about the importance of yard signs versus other factors that are not modeled?
	\vspace{.5cm}
	
	The R2 value indicates how much variation in the dependent variable is explained by our model, the score of which fall between 0 to 1. The closer the score is to one, the higher the variation is explained by the model. Our R2 value of 0.094 indicates a low model fit, indicating that the effect of lawn signs is likely low compared to other factors in explaining the difference in vote share between Cuccinelli and McAuliff.
	
\end{enumerate}  


\end{document}
